\documentclass[border=7bp]{standalone}

\usepackage{fontspec}
\usepackage{xeCJK}
\setmainfont{Calibri}
\setCJKmainfont{SimSun}[AutoFakeBold=1.5]
\usepackage{amsmath}
\usepackage{tikz}
\usetikzlibrary{arrows.meta,calc}

\definecolor{Ink}{HTML}{25313D}
\definecolor{Blue}{HTML}{315E80}
\definecolor{PaleBlue}{HTML}{E7EFF5}
\definecolor{Muted}{HTML}{626B74}
\definecolor{PaleGray}{HTML}{F0F1F3}

\tikzset{
  >=Stealth,
  every node/.style={font=\fontsize{9.5bp}{12.6bp}\selectfont,text=Ink},
  link/.style={->,draw=Blue,dashed,line width=.7pt},
  % 左右两列使用完全相同的宽与高，箭头才会水平
  box/.style={rectangle,draw=Ink,line width=.65pt,align=center,
              text width=4.6cm,minimum height=1.35cm,inner sep=3pt},
  impl/.style={box,draw=Blue,fill=PaleBlue},
  head/.style={rectangle,draw=none,fill=PaleGray,align=center,
               text width=4.6cm,minimum height=0.75cm,inner sep=3pt,
               font=\fontsize{9.5bp}{12.6bp}\selectfont\bfseries},
  note/.style={font=\fontsize{9bp}{12bp}\selectfont,text=Muted},
}

% Shared preamble for the report's editable vector figures.
\usepackage{fontspec}
\usepackage{xeCJK}
\setmainfont{Calibri}
\setCJKmainfont{SimSun}[AutoFakeBold=1.5]
\usepackage{amsmath}
\usepackage{tikz}
\usetikzlibrary{arrows.meta,positioning,calc,shapes.geometric,fit}
\definecolor{Ink}{HTML}{25313D}
\definecolor{Blue}{HTML}{315E80}
\definecolor{PaleBlue}{HTML}{E7EFF5}
\definecolor{Muted}{HTML}{626B74}
\definecolor{PaleGray}{HTML}{F0F1F3}
\tikzset{
  >=Stealth,
  every node/.style={font=\fontsize{10bp}{13bp}\selectfont,text=Ink},
  flow/.style={->,draw=Ink,line width=.65pt},
  frame/.style={rectangle,draw=Ink,line width=.65pt,align=center,
    minimum height=.85cm,text width=5cm,inner sep=5pt},
  emphasis/.style={frame,draw=Blue,fill=PaleBlue},
  supporting/.style={frame,draw=Muted,fill=PaleGray},
  annotation/.style={font=\fontsize{9bp}{12bp}\selectfont,text=Muted,align=center},
}

\begin{document}
\begin{tikzpicture}
  \node[head] at (0,0)   {框架层 \texttt{core}（教师给定）};
  \node[head] at (7.4,0) {实现层 \texttt{stud.g01}};

  % 行距 1.95cm，框高 1.35cm，行间留 0.6cm
  \node[box]  (a1) at (0,-2.00) {\texttt{State}\\状态接口};
  \node[impl] (b1) at (7.4,-2.00) {\texttt{PuzzleBoard}\\不可变棋盘（3.1）};

  \node[box]  (a2) at (0,-3.95) {\texttt{Problem}\\初态、目标、动作、代价};
  \node[impl] (b2) at (7.4,-3.95) {\texttt{NPuzzleProblem}\\\texttt{PuzzleMove}（步价 1，3.4）};

  \node[box]  (a3) at (0,-5.90) {\texttt{Frontier}\\Open 表接口};
  \node[impl] (b3) at (7.4,-5.90) {\texttt{PqFrontier}\\优先队列 + 哈希表（3.2）};

  \node[box]  (a4) at (0,-7.85) {\texttt{Predictor}\\启发估值接口};
  \node[impl] (b4) at (7.4,-7.85) {\texttt{PuzzlePredictor}\\\texttt{LazyPredictor}（3.3）};

  \node[box]  (a5) at (0,-9.80) {\texttt{AbstractSearcher}、\texttt{BestFirstSearcher}\\\texttt{Node}、Closed 表 \texttt{explored}};
  \node[impl] (b5) at (7.4,-9.80) {\texttt{IdAStar}\\阈值迭代 DFS（3.5）};

  \foreach \i in {1,...,5} { \draw[link] (a\i.east) -- (b\i.west); }

  \node[note] at (3.70,-10.85)
    {虚线箭头表示实现层对框架接口的落实，括号内为对应说明小节};
\end{tikzpicture}
\end{document}
